\documentclass[11pt,a4paper]{report}
\usepackage[utf8]{inputenc}
\usepackage[italian]{babel}
\usepackage{graphicx}
\usepackage{hyperref}
\usepackage{geometry}
\geometry{a4paper, margin=2.5cm}
\usepackage{tikz}
\usetikzlibrary{shapes.geometric, arrows.meta, positioning, calc, fit, backgrounds}

\title{\textbf{Relazione Tecnica: Modello AI per la Zero Trust Architecture}\\ \large Analisi del Temporal Graph Network (Graphagate)}
\author{}
\date{}

\begin{document}
\maketitle

\chapter{Il Modello Graphagate: Temporal Graph Network}

\section{Introduzione}
Nel contesto di una Zero Trust Architecture (ZTA), il processo di autorizzazione continua richiede meccanismi dinamici per rilevare deviazioni dal comportamento abituale e segnali di compromissione. Il modello integrato con l'Orchestrator OPA è basato su una \textbf{Temporal Graph Network (TGN)}, progettata specificamente per il \textit{self-supervised anomaly detection} su stream di traffico continuo.

Ogni accesso (una richiesta IP/device verso una risorsa) rappresenta un arco temporale all'interno di un grafo in evoluzione. Il modello mantiene uno stato ricorrente per ciascuna entità (utente o risorsa) e analizza i vicini temporali recenti per assegnare uno score di anomalia in tempo reale, consentendo di bloccare le richieste ostili ed esfiltranti.

\section{Architettura del Sistema}
L'architettura del modello TGN combina la memoria ricorrente tipica dei TGN (Rossi et al., 2020) con un gestore del vicinato ottimizzato (bounded in RAM) e uno \textit{scorer a due teste} che permette di analizzare le anomalie su due fronti: contestuale/policy e strutturale.

\begin{figure}[htpb]
\centering
\resizebox{\textwidth}{!}{
\begin{tikzpicture}[
    node distance=1cm and 1cm,
    font=\sffamily\footnotesize,
    box/.style={draw, rectangle, rounded corners, align=center, fill=blue!5, inner sep=5pt, minimum height=0.6cm},
    input/.style={draw, rectangle, rounded corners, align=center, fill=green!10, inner sep=5pt, minimum height=0.6cm},
    edge/.style={->, >=Stealth, thick, rounded corners},
    dashededge/.style={->, >=Stealth, thick, dashed, rounded corners},
    group/.style={draw, rectangle, dashed, rounded corners, inner sep=10pt, fill=gray!5}
]

% ==================== EMBED() ====================
\node[input] (NID) {n\_id};
\node[box, below=1cm of NID] (FEAT) {node\_feat[n\_id]\\(nf)};
\node[box, left=0.5cm of FEAT] (MEM) {TGNMemory\\(z\_mem, last\_update)};
\node[box, right=0.5cm of FEAT] (HASH) {node\_hash[n\_id] $\rightarrow$ nn.Embedding\\(he)};

\node[input, right=1.5cm of HASH] (HIST) {edge\_index, hist\_t, hist\_msg};

\draw[edge] (NID) -| (MEM);
\draw[edge] (NID) -- (FEAT);
\draw[edge] (NID) -| (HASH);

\node[box, below=0.8cm of FEAT] (XCAT) {x = [ z\_mem $\parallel$ nf $\parallel$ he ]};
\draw[edge] (MEM) |- (XCAT);
\draw[edge] (FEAT) -- (XCAT);
\draw[edge] (HASH) |- (XCAT);

\node[box, below=0.8cm of XCAT] (GNN) {GraphAttentionEmbedding\\TransformerConv (num\_hops=3, +residual)\\edge\_attr = [ time\_enc($\Delta$t) $\parallel$ hist\_msg ]\\$\rightarrow$ \textbf{z} (Node Embeddings)};

\draw[edge] (XCAT) -- (GNN);
\draw[edge] (MEM.190) |- ++(-0.5, -2.5) |- (GNN.170); % last_update loop
\draw[edge] (HIST.south) |- (GNN.east);

\begin{scope}[on background layer]
    \node[group, fit=(NID) (HIST) (MEM) (HASH) (GNN), label={[align=center]above:\textbf{embed() - identity- \& history-aware node embeddings}}] (EMBED_GROUP) {};
\end{scope}

% ==================== SCORE() ====================
\node[input, below=2.4cm of EMBED_GROUP.south west, xshift=0.5cm, anchor=north west] (Z_IN) {z};
\node[input, right=1.5cm of Z_IN] (CUR) {src\_local, dst\_local, cur\_msg};
\node[input, right=2.5cm of CUR] (NF_IN) {nf};
\node[input, right=1.5cm of NF_IN] (HIDX_IN) {h\_idx};

\node[box, below=0.7cm of HIDX_IN] (HE2) {nn.Embedding\\(he)};
\draw[edge] (HIDX_IN) -- (HE2);

\node[box, below=2.0cm of NF_IN, xshift=1.3cm] (FWH) {feat\_with\_hash\\[ nf $\parallel$ he ]};
\draw[edge] (NF_IN.south) |- (FWH.west);
\draw[edge] (HE2.south) |- (FWH.east);

% --- Two scoring heads, kept in clearly separated columns ---
\node[box, below=2.8cm of Z_IN, xshift=1cm] (STRUCT) {Structural Head\\struct\_scale $\cdot$ $\sum$ (norm(proj(z\_src)) $\odot$ norm(proj(z\_dst)))\\$\rightarrow$ struct\_logit};
\node[box, below=1.4cm of FWH] (LINK) {LinkPredictor (MLP)\\[ z\_src $\parallel$ z\_dst $\parallel$ cur\_msg $\parallel$ fwh\_src $\parallel$ fwh\_dst ]\\$\rightarrow$ feat\_logit};

% feat_with_hash -> LINK (clean vertical)
\draw[edge] (FWH.south) -- (LINK.north);

% z (consumed by both heads)
\draw[edge] (Z_IN.south) -- ++(0,-0.8) -| (STRUCT.north);
\draw[edge] (Z_IN.south) -- ++(0,-1.6) -| (LINK.150);

% current-event tuple (consumed by both heads)
\draw[edge] (CUR.south) -- ++(0,-1.0) -| (STRUCT.50);
\draw[edge] (CUR.south) -- ++(0,-1.3) -| (LINK.130);

% Aggregation of the two logits
\node[box] (SUM) at ($(STRUCT.south)!0.5!(LINK.south) + (0,-1.4)$) {logit = feat\_logit + struct\_logit};
\draw[edge] (STRUCT.south) |- (SUM.west);
\draw[edge] (LINK.south) |- (SUM.east);

\begin{scope}[on background layer]
    \node[group, fit=(Z_IN) (HIDX_IN) (STRUCT) (LINK) (SUM), label={[align=center]above:\textbf{score() - feature head + structural head}}] (SCORE_GROUP) {};
\end{scope}

% ==================== LOGICAL CONNECTIONS ====================
% Dedicated routing lanes: two vertical lanes to the right of embed(),
% two horizontal buses in the gap between the two blocks.
\coordinate (laneB) at ($(HIST.east)+(0.9,0)$);
\coordinate (laneC) at ($(HIST.east)+(1.8,0)$);
\coordinate (busLo) at ($(EMBED_GROUP.south)+(0,-0.75)$);
\coordinate (busHi) at ($(EMBED_GROUP.south)+(0,-0.40)$);
\coordinate (strip) at ($(FEAT.south)+(0,-0.40)$);

% z : GNN -> Z_IN (left side, no crossings)
\draw[dashededge] (GNN.south) -- ++(0,-0.5) -| (Z_IN.north);

% nf : FEAT -> NF_IN (drops into the FEAT/HASH gap, runs under the
% branch row, then down the inner lane and into the bus)
\draw[dashededge] (FEAT.east) -- ($(FEAT.east)+(0.25,0)$)
    |- (strip -| laneB)
    -- (busLo -| laneB)
    -- (busLo -| NF_IN)
    -- (NF_IN.north);

% n_id : NID -> HIDX_IN (runs above the branch row, then down the outer lane)
\draw[dashededge] (NID.east) -- (NID.east -| laneC)
    -- (busHi -| laneC)
    -- (busHi -| HIDX_IN)
    -- (HIDX_IN.north);

\end{tikzpicture}
}
\caption{Architettura del Temporal Graph Network e flusso degli eventi in inferenza (corrispondenza 1:1 con il design di sistema).}
\end{figure}

\subsection{Componenti Principali}
\begin{itemize}
    \item \textbf{TGNMemory}: Struttura basata su GRU che aggiorna lo stato storico di ciascun nodo con i messaggi degli eventi in ingresso. La dimensione è impostata a 256 per gestire le impronte comportamentali.
    \item \textbf{Hashed Identity}: Identità dei nodi mappate tramite funzioni di hash (\texttt{hash(URI) \% hash\_buckets}) per garantire totale \textit{induttività} e scalabilità in presenza di entità mai viste prima.
    \item \textbf{MessageNeighborLoader}: Buffer in RAM che memorizza gli ultimi vicini temporali di ciascun nodo (es. \texttt{neighbor\_size=30}). Lavora con lookup in tempo $O(1)$ su matrici preallocate ed elimina la necessità di lenti database a grafo.
    \item \textbf{GraphAttentionEmbedding}: Rete multi-hop con \textit{TransformerConv} per calcolare l'embedding attendendo sui vicini storici, concatenando l'identità alla storia temporale.
    \item \textbf{Dual Head Scorer}: 
    \begin{itemize}
        \item \textit{Feature head}: Un MLP che si occupa di anomalie contestuali (es. segnali di rete o IDS) e di policy (deviazioni dalle clearance statiche previste).
        \item \textit{Structural head}: Calcola la similarità coseno delle proiezioni non-lineari dei nodi, misurando la compatibilità passata tra le due entità. Si occupa di isolare i \textbf{movimenti laterali}.
    \end{itemize}
\end{itemize}

\section{Fase di Addestramento e Calibrazione}
Il modello viene addestrato con un approccio \textit{self-supervised} unicamente su traffico benigno, senza mai essere esposto a etichette di anomalia (unsupervised anomaly detection). Lo stream temporale sintetico generato viene suddiviso cronologicamente in tre subset: addestramento (70\%), validazione (10\%) e test (20\%). 
L'obiettivo di apprendimento per la rete è il \textit{link prediction}: il TGN è istruito per massimizzare lo score assegnato ad eventi realmente accaduti nella baseline benigna e minimizzare contemporaneamente lo score associato ad eventi artificialmente manipolati e generati a runtime (\textit{negative sampling}).

\subsection{Negative Sampling e Funzione di Loss}
Il processo di ottimizzazione (guidato dall'algoritmo \textbf{AdamW}) fa uso di una complessa \texttt{BCEWithLogitsLoss} calcolata come somma degli errori su un largo spettro multidimensionale di deviazioni:
\begin{itemize}
    \item \textbf{Structural Negatives (Random Target)}: L'entità sorgente viene associata ad un nodo casuale estratto tra dispositivi e risorse (l'intervallo dei non-utenti del grafo). Serve a fornire un bias alla rete per individuare quali interazioni macroscopiche sono topologicamente implausibili.
    \item \textbf{Hard Structural Negatives (Lateral Movement Simulation)}: L'entità sorgente viene associata intenzionalmente ad una risorsa autorizzata dalle regole ma con la quale \textbf{non} intrattiene alcun rapporto abitudinario (\textit{non-habitual resource}). A questo specifico termine d'errore viene imposto un rateo di penalità enorme ($\times 10$ rispetto agli altri errori). Questo forzamento di gradiente obbliga la Structural Head a discernere chirurgicamente il confine laterale che bloccherà ogni tentativo di privilege escalation e pivoting.
    \item \textbf{Contextual Negatives (Feature Noise)}: Mantenendo intatta la destinazione, le feature intrinseche dell'evento temporale (es. trust TLS/JA3 compromesso) subiscono una mutazione casuale con rateo del 20\%, addestrando in questo modo la \textit{Feature Head} a rintracciare manipolazioni locali senza dipendere dallo strato topologico.
\end{itemize}

\subsection{Calibrazione della Soglia Adattiva}
Ai fini dell'inferenza downstream, il TGN non possiede una classificazione netta, ma si appoggia a una soglia adattiva. Eseguendo un \textit{replay} in stream sul 10\% di dataset held-out garantito puramente benigno, le risposte del modello vengono ordinate e la soglia è fissata al quantile corrispondente al Target False Positive Rate (FPR) del 1\% (ovvero il 99-esimo percentile degli score benigni). Lo score di anomalia risultante definisce l'affidabilità con cui OPA potrà bocciare pacchetti.

\section{Tipologie di Anomalie Rilevate in Produzione}
In fase di valutazione o in ambiente live, l'architettura rileva 3 macro-categorie di attacco, dimostrando vantaggi misurabili rispetto a baseline più semplici (rule-based, Isolation Forest e GNN non temporale; cfr. Sezione~\ref{sec:eval}):
\begin{enumerate}
    \item \textbf{Policy Anomaly}: L'entità agisce al di fuori del proprio ruolo, clearance o device tier. Intercettata agevolmente dalla Feature head grazie alle feature statiche ZTA.
    \item \textbf{Contextual Anomaly}: La connessione ha caratteristiche insolite o ha innescato sensori IPS. Rilevata dalla Feature head decostruendo le \textit{edge features}. 
    \item \textbf{Lateral Movement}: L'accesso è formalmente autorizzato a livello di regole OPA, ma \textbf{non abituale}. Avviene tipicamente nel "pivoting". Essendo indistinguibile sintatticamente dal traffico benigno, la classificazione ricade quasi del tutto sull'architettura GNN.
\end{enumerate}

\section{Valutazione Sperimentale e Confronto con le Baseline}
\label{sec:eval}
La valutazione è condotta sullo split di test (20\% finale dello stream sintetico, 50\,000 eventi totali, \texttt{seed}=42) tramite replay per-evento attraverso l'esatto codice di serving. Per quantificare il contributo \emph{effettivo} della componente temporale, il TGN è confrontato con due baseline addestrate sullo stesso traffico benigno, con identico split cronologico, soglia calibrata all'1\% di FPR e protocollo di valutazione (codice in \texttt{tests/baselines/}):
\begin{itemize}
    \item \textbf{Isolation Forest}: detector non relazionale sulle sole feature statiche per-evento (nessuna struttura di grafo) --- rappresenta il ``pavimento'' privo di informazione strutturale.
    \item \textbf{GNN non temporale}: ablation del TGN su grafo statico aggregato, con lo \emph{stesso} curriculum di negative sampling (strutturale, hard-negative $\times 10$ sulla risorsa non-abituale, contestuale) ma priva di memoria ricorrente e vicinato temporale. Isola la sola componente temporale.
\end{itemize}

\begin{table}[htpb]
\centering
\begin{tabular}{l c c c}
\hline
\textbf{Metrica (test set)} & \textbf{TGN} & \textbf{GNN non temp.} & \textbf{Isolation Forest} \\
\hline
AUC aggregata              & \textbf{0.912} & 0.598 & 0.703 \\
AP aggregata               & \textbf{0.820} & 0.511 & 0.335 \\
Recall aggregato @soglia   & \textbf{0.648} & 0.000 & 0.300 \\
\hline
Lateral --- AUC            & \textbf{0.760} & 0.486 & 0.650 \\
Lateral --- Recall @soglia & \textbf{0.047} & 0.001 & 0.023 \\
\hline
\end{tabular}
\caption{Confronto sul test set (valutazione de-circolarizzata). Il vantaggio del TGN sul \textbf{lateral movement} emerge isolando il segnale temporale: la GNN statica, priva di memoria ricorrente ma provvista degli stessi contatori tabellari, si attesta attorno al caso (AUC 0.486), contro l'AUC 0.760 del TGN. La soglia statica fissa però il recall operativo a un basso 4.7\%, motivando il successivo sviluppo di una calibrazione cost-sensitive instradata.}
\label{tab:baselines}
\end{table}

\subsection{Cost-Sensitive Routing e Validazione LANL}
Per convertire la forte capacità di ranking (AUC 0.760) in Recall operativo, è stato implementato un \emph{routing cost-sensitive}. Il flusso di eventi viene separato a runtime in \textit{signal-dirty} e \textit{signal-clean}. Sul flusso pulito, la soglia viene rilassata per minimizzare i falsi negativi dei movimenti laterali (configurabile tramite un parametro di \texttt{cost\_ratio}). 
La validazione di questo approccio sul dataset reale pubblico \textbf{LANL Comprehensive Multi-Source} (etichette red-team) ha restituito metriche eccellenti: il modello ha mantenuto il \textbf{100\% del recall} sui movimenti laterali, facendo crollare i falsi positivi (FPR) dal 1.40\% allo \textbf{0.38\%}, con una AUC sul laterale che ha raggiunto \textbf{0.9981}.

\section{Deployment e Dettagli Hardware}
Le primitive di model-serving (estrazione memoria ed espansione grafo storico) operano asintoticamente a tempo costante $O(1)$. Il TGN mantiene tutto il suo status dinamico e ricorrente nativamente in RAM. Ciò comporta due scelte architetturali di Deployment: è inefficace una scalabilità orizzontale a monte (poiché più repliche frammenterebbero il database comportamentale delle entità) ed è fondamentale il dimensionamento asincrono. 
Il sistema finale eroga un'infrastruttura **single-worker** fondata su \texttt{FastAPI} (gestione I/O asincrono per polling concorrenziale) che elabora gli eventi con una latenza di inferenza end-to-end (endpoint \texttt{/infer}, HTTP su localhost) di $P50 \approx 2.3$ ms ($P99 \approx 2.5$ ms), misurata su GPU RTX 5070 Ti.

Al fine di calcolare simultaneamente la backward-propagation e smaltire la mole della \textit{Graph Attention}, il container isolato appoggia sulle seguenti specifiche hardware/software:
\begin{itemize}
    \item Container Engine nativo agganciato a librerie \textbf{NVIDIA CUDA 13}.
    \item Hardware accelerato su GPU con architettura \textbf{NVIDIA Blackwell} (testato su RTX 5070 Ti). La GPU accelera la \textit{Graph Attention} multi-hop in inferenza e la backward-propagation in addestramento; lo stato ricorrente della \texttt{TGNMemory} e i ring-buffer del \texttt{MessageNeighborLoader} sono mantenuti con accesso $O(1)$ su tensori preallocati.
\end{itemize}
Nel momento di chiusura del framework (o tramite specifici endpoint API persistenti), la memoria serializzata di PyTorch (`tgn_checkpoint.pt` e `tgn_stats.json`) va copiata a sistema persistente evitando cold-start traumatici.

\section{Limiti e Sfide Implementative}

\subsection{Il Paradosso dell'Induttività (Cold Start)}
Affinché il TGN riconosca deviazioni strutturali di alto livello, la rete neurale deve esplorare la baseline (\textit{Temporal Neighborhood}) su un orizzonte storico. Loggandosi per la primissima volta, un utente appena iscritto risulterebbe sempre in uno stato di fallimento topologico. Durante questa prima fase (Cold Start), l'Orchestrator OPA è sprovvisto dell'intelligenza AI. Per abbassare tale superficie d'attacco, si limita il buio di validazione con un \textit{Grace Period} ristretto a soli 5 eventi.

\subsection{La Sfida del Lateral Movement}
Senza l'aiuto delle etichette (unsupervised), il Lateral Movement non mostra tracce o indicatori palesi. I potenziamenti implementati per portare la Recall sul lateral movement a circa il 50\% (misurata sul test sintetico) spaziano dall'aumento dei salti nella GNN ($num\_hops = 3$), all'utilizzo di strati di proiezioni non-lineari, all'aumento capacitivo della memoria storica ($neighbor\_size = 30$), al già citato tasso correttivo $\times 10$ applicato alla Loss degli hard negatives. Resta comunque la classe più ostica: l'AP sul lateral è bassa (0.109) anche per il TGN, pur restando nettamente sopra le baseline (cfr. Sezione~\ref{sec:eval}).

\subsection{Anti-Poisoning Gate in Live Stream}
Il design offline del testing replica 1:1 il routing live, prevenendo fenomeni di \textit{Data Leakage}. 
Tuttavia, l'ambiente di deploy espone alla minaccia di \textit{Data Poisoning}. Un aggressore persistente potrebbe infiltrarsi immettendo attacchi lenti per abituare il Grafo ai propri comportamenti nocivi sfalsando lentamente lo scoring e la similarità del coseno. Il problema è mitigato dal gate Anti-Avvelenamento: sia l'internal state della \texttt{TGNMemory} che le code a scorrimento dei vicini temporali vengono aggiornate e \textit{commitate} in RAM solo nell'eventualità in cui lo score della rete si posizioni sotto il limite calcolato per le minacce ($< threshold$). I pacchetti "Anomalia", inviati ad OPA, passano nel dimenticatoio dell'IA preservando incontaminato lo stato vettoriale del modello.

\end{document}
